\begin{tikzpicture}[x=1cm,y=1cm,font=\small]
\draw[gray!45,rounded corners=3pt,fill=white] (0,0) rectangle (13.8,8.2);
\node[anchor=west,font=\small\bfseries] at (.4,7.699999999999999) {顶点情形：方向存在，但不能再用有限斜率表示};

\begin{scope}[shift={(1.6,3.9)},scale=2]
\draw[cncurve,samples=120,domain=-1.2:1.2,variable=\t] plot ({sqrt(2*(1+\t*\t))},\t);
\draw[cntangent] ({sqrt(2)},-1.3)--({sqrt(2)},1.3);
\draw[cnchord] ({9*sqrt(2)/7},-1.25)--({9*sqrt(2)/7},1.25);
\draw[cnsec] ({9*sqrt(2)/7},{4*sqrt(2)/7})--({sqrt(2)},0)--({9*sqrt(2)/7},{-4*sqrt(2)/7});
\fill ({sqrt(2)},0) circle(1.8pt) node[left=6pt] {$P=Q$};
\fill ({9*sqrt(2)/7},{4*sqrt(2)/7}) circle(1.8pt) node[right=5pt] {$M$};
\fill ({9*sqrt(2)/7},{-4*sqrt(2)/7}) circle(1.8pt) node[right=5pt] {$N$};
\end{scope}
\node[anchor=west] at (8.4,6.3) {$P=Q=(\sqrt2,0)$};
\node[anchor=west,orange!85!black] at (8.4,4.9) {$\ell_Q:\ x=\sqrt2$};
\node[anchor=west,red!78!black] at (8.4,3.4) {$MN:\ x=\dfrac{9\sqrt2}{7}$};
\node[anchor=west] at (8.4,1.9) {都竖直，仍然平行};
\node at (6.9,.65) {本图仍取两条割线斜率为 $2,-2$；不写不存在的斜率分式。};

\end{tikzpicture}
